\section{Discussão}\label{discussuxe3o}

\subsection{Desempenho e características dos
dados}\label{desempenho-e-caracteruxedsticas-dos-dados}

Os resultados mostram sinal preditivo, mas com alcance limitado. A
representação híbrida teve o menor RMSLE médio do Ridge na validação
interna; a diferença de 0,004 para \texttt{skill-only} é pequena diante
do desvio-padrão entre dobras (0,028 e 0,031). O experimento sugere
informação nos atributos da skill, mas não demonstra superioridade
robusta do híbrido nem desempenho externo de um modelo textual. A maior
correlação de Spearman de \texttt{repository-only} reforça que as
métricas capturam aspectos diferentes.

A assimetria e a sobredispersão do alvo ajudam a explicar o
comportamento. A maioria das skills reutilizadas aparece em poucos
repositórios, enquanto uma parcela pequena concentra contagens elevadas.
A referência constante teve MAE externo menor que o Random Forest (3,067
contra 3,266), embora RMSLE pior (0,714 contra 0,583). Portanto, o ganho
do modelo depende da métrica escolhida. Acima do P95 do treino, o Random
Forest subestimou todas as 324 observações externas, com MAE de 35,15.

\subsection{Comparação dos modelos e vieses
indutivos}\label{comparauxe7uxe3o-dos-modelos-e-vieses-indutivos}

O benchmark completo foi realizado na visão \texttt{repository-only}. O
Random Forest foi escolhido como candidato principal pelo menor RMSLE da
validação agrupada (0,5753). No teste externo, obteve 0,5833; o Gradient
Boosting obteve 0,5907, mas MAE menor (3,1820 contra 3,2663). Esses
valores externos avaliam candidatos já congelados e não são uma nova
rodada de seleção. A variação entre três dobras e a única divisão
externa limitam a precisão de comparações pequenas.

O Ridge impõe uma relação aditiva e regularizada, sendo estável e
relativamente fácil de interpretar, mas tende a encolher valores
extremos. KNN depende da geometria local dos atributos e teve o pior
RMSLE do benchmark completo. A árvore de decisão isolada não superou os
ensembles e apresentou maior variabilidade. Random Forest e Gradient
Boosting capturam não linearidades e interações, mas podem explorar
sinais contextuais que não generalizam para outros ecossistemas.

\subsection{Explicabilidade e
interpretação}\label{explicabilidade-e-interpretauxe7uxe3o}

As análises textuais exploratórias identificaram termos editoriais e
sinais de templates; a robustez ao mascaramento não mostra que sejam
propriedades causais da skill. No Random Forest
\texttt{repository-only}, a importância por impureza destacou
\texttt{repo\_stars}, \texttt{history\_fetched}, ausência de
\texttt{commit\_count} e \texttt{repo\_forks}. Já a permutação agrupada
colocou \texttt{repo\_language} em primeiro (aumento médio de RMSLE
0,029), seguido por \texttt{commit\_count} (0,014),
\texttt{history\_fetched} (0,013) e \texttt{repo\_stars} (0,012). A
divergência decorre em parte das propriedades das medidas de importância
e da correlação entre atributos; o sinal de estrelas variou
substancialmente entre dobras.

No cenário que varia estrelas mantendo as outras entradas, a previsão
média foi não monotônica. Isso enfraquece uma regra simples sobre
popularidade e reutilização, embora o cenário não seja uma estimativa
causal e possa criar combinações raras de atributos. O uso de
\texttt{repo\_language} e histórico também pode codificar origem,
cobertura da coleta ou práticas de distribuição. A interpretação do
Random Forest é apenas contextual; não identifica variáveis textuais ou
estruturais da skill como determinantes da reutilização.

\subsection{Implicações para a engenharia de
software}\label{implicauxe7uxf5es-para-a-engenharia-de-software}

O pipeline pode apoiar análises exploratórias de autoria e descoberta,
indicando convenções que merecem inspeção ou priorizando skills para
auditoria manual. Ele ainda não justifica um linter automático, um
ranking de qualidade ou uma intervenção de autoria. Essas aplicações
exigiriam validação prospectiva, avaliação de utilidade e controle de
possíveis efeitos adversos, como reforçar padrões de popularidade ou
favorecer determinados ecossistemas.

Para o trabalho, a contribuição principal é metodológica: comparar
modelos sob a mesma divisão agrupada, relacionar desempenho a vieses
indutivos e demonstrar que métricas agregadas podem ocultar erros graves
na cauda. Os resultados também delimitam o que pode ser inferido de uma
medida observacional de replicação exata.

\subsection{Ameaças à validade e
limitações}\label{ameauxe7as-uxe0-validade-e-limitauxe7uxf5es}

\texttt{repository\_count} mede a presença de conteúdo idêntico em
repositórios distintos. Não mede intenção de adoção, uso efetivo,
qualidade, utilidade, sucesso de execução ou transferência de
conhecimento. Cópias modificadas e equivalentes semânticos não são
capturados.

O agrupamento por \texttt{representative\_repo\_id} elimina a
sobreposição desse identificador entre treino e validação, mas não
garante que todos os repositórios relacionados, organizações ou cadeias
de cópia estejam separados. Metadados de repositório também podem
funcionar como proxies de popularidade ou origem.

O RMSLE reduz a influência de valores extremos, mas pode ocultar erros
absolutos elevados. Por isso, a análise também utiliza MAE, Spearman,
uma referência constante e auditoria por quantis. A subestimação acima
do P95 mostra que nenhuma métrica isolada caracteriza adequadamente o
desempenho. A comparação textual foi avaliada apenas com Ridge dentro do
desenvolvimento; não autoriza afirmar que o modelo final explica
características intrínsecas da skill.

Os dados provêm de um snapshot específico do GitHub e de um ecossistema
particular de Agent Skills. A generalização para outros períodos, hosts,
comunidades ou mecanismos de distribuição não é garantida. A avaliação
externa reportada é válida para a divisão adotada neste estudo e não
substitui replicações em outros snapshots ou ecossistemas.
